% Maschinenlesbarer Quellenkorpus fuer Abgabeliste und Seitenbudget.
% Diese Datei ist die Single Source of Truth fuer alle gezaehlten Quellen.
%
% Format:
%   \AbgabeQuelle{Frage}{Status}{BibTeX-Key}{Seitenbereich}{Seitenzahl}
%
% Status:
%   K = Kernliteratur, gezaehlt
%   S = gezaehlte Stuetzliteratur
%   N = nicht gezaehlt / nur Hinweis
%
% Die Befehle werden in mpv.tex als No-op definiert. Die Auswertung erfolgt
% ausschliesslich ueber abgabe_korpus.py.

% Frage 1: Multiperspektivische Erfassung verdeckter Begabungen
\AbgabeQuelle{F1}{K}{stamm2021fehlenderblick}{576--585}{10}%
\AbgabeQuelle{F1}{K}{stamm2025vonuntennachoben}{36--37, 58--62}{7}%
\AbgabeQuelle{F1}{K}{preckel2013hochbegabung}{15--47}{33}%
\AbgabeQuelle{F1}{K}{gauckreimann2021psychdiagnostik}{239--245}{7}%
\AbgabeQuelle{F1}{K}{haag2018leistungsstanddiagnostik}{153--165, 187--188}{15}%
\AbgabeQuelle{F1}{K}{kellerkoller2025hellekoepfe}{76--78}{3}%
\AbgabeQuelle{F1}{K}{baumschader2021twice}{588--598}{11}%
\AbgabeQuelle{F1}{K}{koop2025herkunft}{64--67}{4}%
\AbgabeQuelle{F1}{K}{stern2025intelligenz}{15--18}{4}%
\AbgabeQuelle{F1}{K}{webb2020doppeldiagnosen}{87--94}{8}%
\AbgabeQuelle{F1}{K}{warneckehauke2020bildungsgerechtigkeit}{242--253}{12}%
\AbgabeQuelle{F1}{K}{kellerkoller2013erkennen}{1--2}{2}%
\AbgabeQuelle{F1}{K}{muelleroppliger2021paeddiagnostik}{224--235}{12}%
\AbgabeQuelle{F1}{N}{baudson2021wasdenken}{115--128}{0}%
\AbgabeQuelle{F1}{N}{muelleroppliger2021begabungsmodelle}{204--219}{0}%
\AbgabeQuelle{F1}{N}{baudson2025besserfinden}{35--40}{0}%
\AbgabeQuelle{F1}{N}{trautmann2016einfuehrung}{34--41}{0}%
\AbgabeQuelle{F1}{N}{lehwald2017motivation}{78--92}{0}%
\AbgabeQuelle{F1}{K}{kappus2010migration}{63--70, 74}{9}%
\AbgabeQuelle{F1}{N}{grossenbacher2014integrative}{317--325}{0}%
\AbgabeQuelle{F1}{N}{erzinger2023pisa}{45--48}{0}%
\AbgabeQuelle{F1}{N}{bfs2022migration}{34--35}{0}%
\AbgabeQuelle{F1}{N}{dvs2025bbf}{online}{0}%
\AbgabeQuelle{F1}{N}{maehler2018diagnostik}{153--165,187f.}{0}%

% Frage 2: Sensomotorische und schriftsprachliche Barrieren
\AbgabeQuelle{F2}{K}{saegesserwyss2021grafinkrahmenmodell}{2--3, 6--11}{8}
\AbgabeQuelle{F2}{K}{nottbusch2017graphomotorik}{128--136}{9}
\AbgabeQuelle{F2}{K}{hurschler2020handschriftbeurteilung}{1--24}{24}
\AbgabeQuelle{F2}{K}{gold2018lesenkannmanlernen}{50--53, 62, 67--88}{27}
\AbgabeQuelle{F2}{K}{lehwald2017motivation}{70--72, 84--89}{9}
\AbgabeQuelle{F2}{K}{greiten2021underachievement}{546--553}{8}
\AbgabeQuelle{F2}{N}{baumschader2021twice}{588--598}{0}
\AbgabeQuelle{F2}{N}{warneckehauke2020bildungsgerechtigkeit}{247--252}{0}
\AbgabeQuelle{F2}{N}{hoyer2013begabung}{111--113}{0}
\AbgabeQuelle{F2}{N}{sturm2016graphomotorik}{183--198}{0}

% Frage 3: Relationale Bedingungen von Potenzialsichtbarkeit
\AbgabeQuelle{F3}{K}{grossrieder2010anerkennung}{88--94}{7}
\AbgabeQuelle{F3}{K}{baudson2021wasdenken}{115--129}{15}
\AbgabeQuelle{F3}{K}{kuhl2021begabungbildungbeziehung}{185--201}{17}
\AbgabeQuelle{F3}{K}{wagener2021bfhemmendfoerdernd}{418--424}{7}
\AbgabeQuelle{F3}{K}{behrensen2019inklusive}{86--98}{13}
\AbgabeQuelle{F3}{K}{kuhl2019diversitaet}{35--54}{20}
\AbgabeQuelle{F3}{K}{boosnuenning2022interethnisch}{51--65}{15}
\AbgabeQuelle{F3}{K}{kesselshannover2015gleichaltrige}{288}{1}
\AbgabeQuelle{F3}{K}{tschoppbuholzergruetter2022intergruppenkontakt}{35--40}{6}
\AbgabeQuelle{F3}{N}{baumschader2021twice}{588--598}{0}
\AbgabeQuelle{F3}{N}{kappus2010migration}{63--70, 74}{0}
\AbgabeQuelle{F3}{K}{baumert2022freundschaftwerte}{40--49}{10}
\AbgabeQuelle{F3}{N}{gysinscherzinger2022freundschaftenwertvoll}{8--12}{0}
\AbgabeQuelle{F3}{K}{evers2025stress}{21--27}{7}
\AbgabeQuelle{F3}{K}{lehwald2017motivation}{57--63}{7}

% Frage 4: Inklusive Begabungsfoerderung am Beispiel SOLUX
\AbgabeQuelle{F4}{K}{buholzerkummerwyss2010reaktionen}{78--85}{8}
\AbgabeQuelle{F4}{K}{muelleroppliger2021plurale}{32--42}{11}
\AbgabeQuelle{F4}{K}{muelleroppliger2021begabungsmodelle}{204--219}{16}
\AbgabeQuelle{F4}{K}{reisrenzullimueller2021sem}{333--345}{13}
\AbgabeQuelle{F4}{K}{muelleroppliger2021paeddiagnostik}{224--235}{0}
\AbgabeQuelle{F4}{K}{weigand2021separativ}{290--298}{9}
\AbgabeQuelle{F4}{K}{muelleroppliger2021adaptive}{374--385}{12}
\AbgabeQuelle{F4}{K}{schulteterhardt2020potenzialentwicklung}{264--267}{4}
\AbgabeQuelle{F4}{K}{lehwald2017motivation}{151--158}{8}
\AbgabeQuelle{F4}{K}{hoyer2013begabung}{94--99}{6}
\AbgabeQuelle{F4}{K}{sedmak2021bildungsgerechtigkeit}{65--75}{11}
\AbgabeQuelle{F4}{K}{grossenbacher2014integrative}{317--325}{9}
\AbgabeQuelle{F4}{N}{horvath2021elite}{77--87}{0}
\AbgabeQuelle{F4}{N}{weigand2021person}{46--59}{0}
\AbgabeQuelle{F4}{N}{baudson2021wasdenken}{119}{0}
\AbgabeQuelle{F4}{N}{dvs2025bbf}{online}{0}


% Frage 5: Beratung und Zusammenarbeit aus SHP-Perspektive
\AbgabeQuelle{F5}{K}{weigand2021person}{46--59}{14}
\AbgabeQuelle{F5}{K}{horvath2021elite}{77--87}{11}
\AbgabeQuelle{F5}{K}{muellerboeschschaffnermenn2021udl}{94--116}{23}
\AbgabeQuelle{F5}{K}{macha2019gender}{160--171}{12}
\AbgabeQuelle{F5}{K}{groschefussangelgraesel2020kokonstruktion}{463--467, 469--472}{9}
\AbgabeQuelle{F5}{K}{nguyensliwka2021massnahmen}{348--356}{9}
\AbgabeQuelle{F5}{K}{widmerwolf2018multiprofessionell}{299--303, 308--309}{7}
\AbgabeQuelle{F5}{K}{kosoroklabhart2021voneltern}{14--15, 38--39}{4}
\AbgabeQuelle{F5}{K}{baudson2025besserfinden}{35--40}{6}
\AbgabeQuelle{F5}{N}{muelleroppliger2021paeddiagnostik}{224--238}{0}
\AbgabeQuelle{F5}{N}{reisrenzullimueller2021sem}{333--347}{0}
\AbgabeQuelle{F5}{N}{weigand2021separativ}{290--298}{0}
\AbgabeQuelle{F5}{N}{muelleroppliger2021plurale}{32--45}{0}
\AbgabeQuelle{F5}{N}{sedmak2021bildungsgerechtigkeit}{65--75}{0}
\AbgabeQuelle{F5}{N}{muelleroppliger2021adaptive}{384--385}{0}
\AbgabeQuelle{F5}{N}{wagener2021bfhemmendfoerdernd}{418--424}{0}
\AbgabeQuelle{F5}{N}{kummerwyss2017kooperativunterrichten}{151--154, 156--160}{0}
